\begin{figure}[t]
  \centering
  \begin{tikzpicture}[
    node distance=2.0cm,
    font=\small,
    >=stealth,
    line width=1.0pt,
  ]
    % Define colors (RGB converted from hex)
    \colorlet{garnet}{Garnet}
    \colorlet{atlantic}{Atlantic}
    \colorlet{rose}{Rose}
    \definecolor{black70}{RGB}{54, 54, 54}

    % Main circular flow positions
    % Using polar coordinates for circular layout
    \coordinate (center) at (0, 0);

    % Node 1: Radar Detections (top)
    \node[
      draw=garnet,
      fill=white,
      rectangle,
      line width=1.5pt,
      minimum width=2.5cm,
      minimum height=0.8cm,
      anchor=center
    ] (radar) at (0, 3.5) {Radar Detections};

    % Node 2: Cross-Modal Projection (right-top)
    \node[
      draw=black70,
      fill=white,
      rectangle,
      line width=1.5pt,
      minimum width=2.8cm,
      minimum height=0.8cm,
      anchor=center
    ] (projection) at (3.2, 2.2) {Cross-Modal Projection};

    % Node 3: Noisy Bounding Boxes (right)
    \node[
      draw=garnet,
      fill=white,
      rectangle,
      line width=1.5pt,
      minimum width=2.8cm,
      minimum height=0.8cm,
      anchor=center
    ] (noisyboxes) at (3.8, 0.5) {Noisy Bounding Boxes};

    % Node 4: Train Camera Detector (right-bottom)
    \node[
      draw=atlantic,
      fill=white,
      rectangle,
      line width=1.5pt,
      minimum width=3.0cm,
      minimum height=0.8cm,
      anchor=center
    ] (train) at (2.5, -2.0) {Train Camera Detector (Student)};

    % Node 5: Generate Pseudo-Labels (bottom-left)
    \node[
      draw=atlantic,
      fill=white,
      rectangle,
      line width=1.5pt,
      minimum width=3.2cm,
      minimum height=0.8cm,
      anchor=center
    ] (pseudo) at (-2.5, -2.0) {Generate Pseudo-Labels};

    % Confidence threshold annotation for pseudo-labels
    \node[
      draw=rose,
      fill=white,
      rectangle,
      line width=1.0pt,
      minimum width=1.8cm,
      minimum height=0.5cm,
      font=\footnotesize,
      anchor=south
    ] (tau) at (-2.5, -2.8) {$\tau$ threshold};

    % Node 6: Combine Labels (left)
    \node[
      draw=black70,
      fill=white,
      rectangle,
      line width=1.5pt,
      minimum width=2.8cm,
      minimum height=0.8cm,
      anchor=center
    ] (combine) at (-3.5, 0.5) {Pseudo-Labels + Radar Labels};

    % Teacher model (side branch)
    \node[
      draw=atlantic,
      fill=white,
      rectangle,
      line width=1.5pt,
      minimum width=2.5cm,
      minimum height=0.8cm,
      anchor=center
    ] (teacher) at (0, -3.8) {Teacher Model};

    % EMA Update label
    \node[
      draw=atlantic,
      fill=white,
      rectangle,
      line width=1.0pt,
      minimum width=2.0cm,
      minimum height=0.5cm,
      font=\footnotesize,
      anchor=south
    ] (ema) at (2.0, -2.9) {EMA Update};

    % Main circular flow arrows
    \draw[->,
      line width=1.5pt,
      garnet
    ] (radar) -- (projection);

    \draw[->,
      line width=1.5pt,
      garnet
    ] (projection) -- (noisyboxes);

    \draw[->,
      line width=1.5pt,
      atlantic
    ] (noisyboxes) -- (train);

    \draw[->,
      line width=1.5pt,
      atlantic
    ] (train) -- (pseudo);

    \draw[->,
      line width=1.5pt,
      black70
    ] (pseudo) -- (combine);

    \draw[->,
      line width=1.5pt,
      garnet
    ] (combine) -- (radar);

    % Teacher-Student connection
    \draw[->,
      line width=1.0pt,
      atlantic,
      dashed
    ] (train) -- (teacher);

    % Teacher to pseudo-labels
    \draw[->,
      line width=1.0pt,
      atlantic,
      dashed
    ] (teacher) -- (pseudo);

    % Iteration annotations
    \node[
      font=\footnotesize,
      black70,
      anchor=center
    ] (iter1) at (-4.8, 2.5) {Iteration 1:};

    \node[
      font=\footnotesize,
      black70,
      anchor=center
    ] (iter1text) at (-4.8, 2.0) {Radar labels only};

    \node[
      font=\footnotesize,
      black70,
      anchor=center
    ] (iter2) at (-4.8, -1.0) {Iteration 2+:};

    \node[
      font=\footnotesize,
      black70,
      anchor=center
    ] (iter2text) at (-4.8, -1.5) {Radar +};

    \node[
      font=\footnotesize,
      black70,
      anchor=center
    ] (iter2text2) at (-4.8, -1.9) {pseudo-labels};

    % Risk annotation box
    \node[
      draw=rose,
      fill=white,
      rectangle,
      line width=1.5pt,
      minimum width=3.0cm,
      minimum height=1.2cm,
      font=\footnotesize,
      align=center,
      anchor=center
    ] (risk) at (5.8, -0.8) {Risk: Systematic Bias Amplification};

    % Arrow pointing from noisyboxes to risk
    \draw[->,
      line width=1.0pt,
      rose,
      dashed
    ] (noisyboxes) -- (risk);

  \end{tikzpicture}
  \caption{Iterative self-training loop for radar-guided camera labeling. The teacher model, updated via exponential moving average (EMA) of the student, generates pseudo-labels that supplement noisy radar-projected annotations. Systematic bias from calibration errors risks amplification across iterations.}
  \label{fig:self_training_loop}
\end{figure}
